\chapter{Composizione dell'ambiente}
\label{cap:ambiente}

A partire dai dati descritti nel Capitolo~\ref{cap:dati}, questo capitolo illustra
l'\textbf{ambiente di simulazione} vero e proprio: l'interfaccia software che trasforma i
campi idrodinamici in un problema di controllo sequenziale, la struttura della squadra di
agenti, lo spazio delle azioni, i vincoli di navigazione, la procedura di inizializzazione
e il segnale di reward. Si tratta dell'infrastruttura \textbf{condivisa} dai due approcci
confrontati nella tesi: sia il metodo del gradiente (Capitolo~\ref{cap:fcm}) sia l'agente
PPO (Capitolo~\ref{cap:rl}) operano su questo stesso ambiente, con la stessa dinamica e
gli stessi vincoli. L'unica differenza risiede nel modulo che, ricevuta l'osservazione,
sceglie l'azione: una regola analitica nel primo caso, una rete neurale addestrata nel
secondo.

\section{L'ambiente software}
\label{sec:env}

L'ambiente \`e implementato dalla classe \texttt{SourceSeekingEnv} secondo l'interfaccia
standard \texttt{gymnasium.Env}. Esso riceve in ingresso un \texttt{ConcentrationField}
(con i relativi \texttt{WindData} e \texttt{CurrentData}), oppure un \texttt{DataManager}
da cui estrarre dinamicamente un campo casuale ad ogni episodio --- modalit\`a usata in
addestramento per variare di continuo sorgente e scenario. Il dominio \`e la griglia
regolare $300 \times 250$ celle a passo $10\,\text{m}$ ($3 \times 2{,}5\,\text{km}$)
descritta nella Sezione~\ref{sec:sorgenti}.

L'oggetto ambiente viene quindi incapsulato in una catena di \emph{wrapper}:
\begin{itemize}
    \item un \texttt{ActionMasker} che applica l'action masking
          (Sezione~\ref{sec:masking}), escludendo le azioni non ammissibili;
    \item un \emph{vettorizzatore} (\texttt{DummyVecEnv}) che uniforma l'interfaccia a
          quella attesa dalla libreria di RL;
    \item per il solo addestramento RL, un \texttt{VecNormalize} che normalizza le
          osservazioni con media e varianza stimate online.
\end{itemize}

\begin{lstlisting}
raw = SourceSeekingEnv(config, concentration_field=field,
                       wind_data=wind, current_data=current)
env = ActionMasker(raw, mask_fn)               # masking delle azioni
env = DummyVecEnv([lambda: env])               # vettorizzazione
env = VecNormalize(env, norm_reward=False)     # solo per il training RL
\end{lstlisting}

\subsection{Architettura master--slave}
\label{sec:master_slave}

Il sistema \`e organizzato come una \textbf{squadra di agenti} in formazione
\emph{master--slave}. L'agente \textbf{master} \`e l'unico a decidere il movimento e a
spostarsi attivamente sul dominio; attorno ad esso, otto agenti \textbf{slave}
mantengono una formazione rigida (Figura~\ref{fig:formazione}), posizionandosi a distanza
$r_s$ (il \emph{sensor range}) lungo le otto direzioni cardinali e diagonali $\hat{d}_i$ (N, S, E, O, NE, SE,
NO, SO). Ciascuno slave campiona la concentrazione del contaminante nella propria
posizione $\mathbf{p} + r_s\,\hat{d}_i$ e la riporta al master, senza alcuna autonomia
decisionale propria: il loro unico compito \`e fornire le misurazioni e conservare la
formazione rispetto al master mentre questo si muove.

Le otto misurazioni degli slave, insieme alla lettura del master nella posizione
centrale, costituiscono l'informazione locale sul campo di concentrazione su cui si
basano entrambi gli approcci. Nel metodo del gradiente tali letture vengono usate per
stimare il gradiente locale; nell'approccio RL confluiscono nello spazio di osservazione
della rete neurale che governa il master. Le posizioni degli slave e il relativo
campionamento si esprimono come:

\begin{lstlisting}
DIAG = 1 / np.sqrt(2)
DIRECTIONS = np.array([(0,1),(0,-1),(1,0),(-1,0),          # N, S, E, O
                       (DIAG,DIAG),(DIAG,-DIAG),
                       (-DIAG,DIAG),(-DIAG,-DIAG)])         # diagonali
# ogni slave campiona a distanza r_s lungo la propria direzione
readings = [field.get_concentration(*(master_pos + r_s * d)) for d in DIRECTIONS]
\end{lstlisting}

\begin{figure}[H]
\centering
\begin{tikzpicture}[
  font=\footnotesize,
  slave/.style={circle, draw, fill=blue!14, minimum size=6.5mm, inner sep=0pt},
  master/.style={circle, draw, fill=orange!75, minimum size=9mm, inner sep=0pt},
  link/.style={-, gray, densely dashed}
]
  \def\R{2.3}
  \draw[gray!45, densely dotted] (0,0) circle (\R);
  \foreach \a/\lab in {90/N, 270/S, 0/E, 180/O, 45/NE, 315/SE, 135/NO, 225/SO} {
    \draw[link] (0,0) -- (\a:\R);
    \node[slave] at (\a:\R) {\scriptsize S};
    \node[gray!60!black] at (\a:{\R+0.42}) {\lab};
  }
  \node[master] (m) at (0,0) {\scriptsize M};
  \draw[{Latex}-{Latex}, thick] (0,0) -- (20:\R)
       node[midway, above, sloped, fill=white, inner sep=1pt]{$r_s$};
\end{tikzpicture}
\caption{Formazione master--slave. L'agente \textbf{master} (M, arancio) occupa il centro;
         otto agenti \textbf{slave} (S, blu) sono disposti a distanza $r_s$ (il \emph{sensor
         range}) lungo le direzioni cardinali e diagonali. Ogni slave campiona la
         concentrazione nella propria posizione e la riporta al master; le stesse otto
         direzioni sono anche le azioni di movimento disponibili.}
\label{fig:formazione}
\end{figure}

Questa unica corona di 8 sensori \`e la configurazione base (\emph{singola corona}); il
Capitolo~\ref{cap:rl} valuta anche una variante a \textbf{doppia corona} che aggiunge una
seconda formazione di 8 slave a $r_{s2} = 50\,\text{m}$, per dare al master informazione
sul plume a scala spaziale pi\`u larga.

\subsection{Spazio delle azioni e dinamica dell'episodio}
\label{sec:dinamica}

Il master si muove con uno spazio delle azioni discreto a \textbf{8 direzioni}
(le stesse $\hat{d}_i$ della formazione), con velocit\`a nominale di $1\,\text{m/s}$ e
passo temporale $\Delta t = 10\,\text{s}$, spostandosi quindi di $10\,\text{m}$ per
azione (o $\sim$7\,m per componente nelle direzioni diagonali). Ogni episodio ha durata
massima di 1080 step ($= 3\,\text{h}$ di tempo simulato) e termina con \textbf{successo}
quando il master raggiunge la sorgente (distanza $\leq 50\,\text{m}$). Il campo di
concentrazione evolve durante l'episodio: ogni alcuni step dell'agente il campo MIKE21
avanza di un frame temporale, mantenendo vento e correnti sincronizzati. La conversione
azione\,$\to$\,spostamento \`e immediata:

\begin{lstlisting}
ACTION_MAP = {0:(0,1), 1:(0,-1), 2:(1,0), 3:(-1,0),       # N, S, E, O
              4:(DIAG,DIAG), 5:(DIAG,-DIAG),
              6:(-DIAG,DIAG), 7:(-DIAG,-DIAG)}            # diagonali
dx, dy = ACTION_MAP[action]
step = max_velocity * dt                  # 1 m/s * 10 s = 10 m
master_pos += np.array([dx, dy]) * step
\end{lstlisting}

Questo spazio d'azione a passo fisso \`e il punto di partenza; il Capitolo~\ref{cap:rl} lo
estende a \textbf{velocit\`a variabile} rendendo la velocit\`a stessa una scelta
dell'agente (spazio $\text{Discrete}(8\times K)$), cos\`i che il master decida a ogni passo
non solo la direzione ma anche l'entit\`a dello spostamento.

\subsection{Action masking}
\label{sec:masking}

Per impedire le collisioni con la costa, l'ambiente calcola ad ogni passo una
\emph{maschera} delle azioni: le direzioni che porterebbero il master sulla terraferma o
fuori dal dominio vengono escluse dall'insieme delle azioni ammissibili. Questo
meccanismo, sfruttato sia dal metodo del gradiente sia dall'agente RL, garantisce che
vengano percorse soltanto traiettorie valide senza dover penalizzare a posteriori le
azioni non consentite.

\begin{lstlisting}
def action_masks(self):
    masks = np.ones(8, dtype=bool)
    step = self.max_velocity * self.dt
    for a, (dx, dy) in ACTION_MAP.items():
        nx, ny = self.x + dx*step, self.y + dy*step
        if out_of_domain(nx, ny) or self.field.is_land(nx, ny):
            masks[a] = False                  # azione invalida
    return masks
\end{lstlisting}

La maschera \`e costruita a partire dalla \emph{land mask} del campo e dalla mappa delle
distanze dalla costa precalcolata via \emph{Euclidean Distance Transform} (EDT,
Sezione~\ref{sec:dataloader}): la stessa mappa fornisce anche la penalit\`a di prossimit\`a
alla terra usata nel segnale di reward (Sezione~\ref{sec:reward}).

\subsection{Inizializzazione dell'agente}
\label{sec:spawn}

Il master viene inizializzato su una cella appartenente al plume (concentrazione
$c \geq 0{,}5$), a distanza dalla costa $\geq 50\,\text{m}$ e a una distanza dalla
sorgente controllata. Nel corso del progetto si sono succeduti \textbf{due algoritmi di
spawn}, e quale dei due sia in uso dipende dalla fase di lavoro: gli addestramenti iniziali
(a velocit\`a fissa) adottano una \emph{cascata di anelli} a soglie fisse, mentre a partire
dall'analisi a velocit\`a variabile (Sezione~\ref{sec:rl_velocity}) si passa a un
\emph{disco adattivo} $[d_{\min},d_{\max}]$ con campionamento uniforme in distanza. Vanno
descritti entrambi perch\'e coesistono nei risultati riportati: ogni modello eredita
l'algoritmo in vigore quando \`e stato addestrato, e i confronti fra modelli tengono conto
di questa differenza.

\paragraph{Primo algoritmo (addestramenti a velocit\`a fissa): cascata di anelli.}
La procedura iniziale (Figura~\ref{fig:spawn_cascade}) scandisce una lista di anelli in
ordine decrescente di distanza dalla sorgente, restituendo una cella a caso del
\emph{primo anello non vuoto}:
\begin{center}
$(2000\text{--}2500) \to (1500\text{--}2000) \to (1000\text{--}1500) \to
(500\text{--}1000) \to (250\text{--}500) \to (50\text{--}250) \to (0\text{--}50)\;\text{m}$
\end{center}
Per ogni anello vengono selezionate le celle con $c \geq 0{,}5$ e distanza dalla terra
$\geq 50\,\text{m}$; la cascata scende all'anello successivo solo se nessuna cella del
corrente soddisfa entrambe le condizioni. La distanza iniziale media di spawn risulta di
circa $800\,\text{m}$.

\begin{lstlisting}
CASCADE = [(2000,2500),(1500,2000),(1000,1500),
           (500,1000),(250,500),(50,250),(0,50)]
for dmin, dmax in CASCADE:                    # dal piu' lontano al piu' vicino
    cells = plume_cells[(dist >= dmin) & (dist <= dmax) & (land_dist >= 50)]
    if len(cells):
        return random.choice(cells)           # primo anello non vuoto
\end{lstlisting}

\begin{figure}[H]
\centering
\begin{tikzpicture}[font=\footnotesize]
  % anelli concentrici (1 unita' = 500 m); disegnati dal piu' grande al piu' piccolo
  \foreach \r/\col in {5/blue!22, 4/blue!18, 3/blue!14, 2/blue!11, 1/blue!8, 0.5/blue!5}{
    \fill[\col] (0,0) circle (\r);
  }
  \foreach \r in {0.1,0.5,1,2,3,4,5}{ \draw[gray!55] (0,0) circle (\r); }
  % sorgente
  \node[star, star points=5, fill=yellow, draw, minimum size=5mm, inner sep=0pt] at (0,0) {};
  % etichette distanze (m) sull'asse
  \foreach \r/\d in {0.5/250, 1/500, 2/1000, 3/1500, 4/2000, 5/2500}{
    \node[fill=white, inner sep=0.8pt, font=\scriptsize] at (\r,0) {\d};
  }
  % freccia: ordine di selezione dal piu' lontano al piu' vicino
  \draw[-{Latex[length=2.4mm]}, thick, orange!85!black] (4.6,0.55) to[bend right=12] (0.9,0.55);
  \node[orange!70!black, align=center, font=\scriptsize] at (2.7,1.35)
       {ordine di selezione\\(dal pi\`u lontano)};
\end{tikzpicture}
\caption{Cascata di anelli concentrici per l'inizializzazione. Il master \`e collocato in
         un anello attorno alla sorgente ($\star$), partendo dal pi\`u lontano
         ($2000$--$2500\,\text{m}$) e scendendo verso i pi\`u vicini solo se nessuna cella
         dell'anello corrente soddisfa i vincoli (concentrazione $\geq 0{,}5$, distanza dalla
         terra $\geq 50\,\text{m}$). Le distanze sono in metri dalla sorgente.}
\label{fig:spawn_cascade}
\end{figure}

\textbf{Pro.} Lo schema \`e semplice e deterministico nell'ordinamento, e privilegia in modo
naturale gli spawn lontani: con plume esteso l'agente parte quasi sempre dal bordo esterno,
affrontando il compito nella sua versione pi\`u difficile e informativa.
\textbf{Contro.} Restituendo sempre il \emph{primo} anello popolato, gli spawn si
\textbf{addensano quasi tutti alla stessa distanza} (l'anello esterno del plume del
momento): la distribuzione della distanza iniziale \`e di fatto degenere e copre male
l'intervallo intermedio. Le soglie sono inoltre \emph{fisse}, e non si adattano a plume
pi\`u compatti o pi\`u diffusi del previsto.

\paragraph{Secondo algoritmo (dall'analisi a velocit\`a variabile): disco adattivo $[d_{\min},\,d_{\max}]$.}
Per distribuire i punti di partenza su tutto l'intervallo utile, la versione definitiva
costruisce un \emph{disco di spawn} $[d_{\min},\,d_{\max}]$ che si adatta allo scenario e
campiona in modo \textbf{uniforme rispetto alla distanza} (Figura~\ref{fig:spawn_disk}).
Detto $\{d_i\}$ l'insieme delle distanze dalla sorgente delle celle di plume ammissibili
($c \geq 0{,}5$, distanza dalla terra $\geq 50\,\text{m}$):
\begin{itemize}
  \item $d_{\max} = \min\!\big(3000\,\text{m},\ \max_i d_i\big)$: il tetto di $3\,\text{km}$,
        abbassato all'estensione massima del plume quando questo non arriva cos\`i lontano;
  \item $d_{\min} = \max\!\big(300\,\text{m},\ \mathrm{perc}_{50}(\{d_i\})\big)$: la mediana
        delle distanze delle celle di plume, con un \emph{floor} assoluto a $300\,\text{m}$.
        Cos\`i $d_{\min}$ \textbf{scala con la diffusione del plume}: plume molto esteso
        $\Rightarrow$ celle lontane $\Rightarrow$ mediana alta $\Rightarrow$ $d_{\min}$ alta;
        plume compatto $\Rightarrow$ $d_{\min}$ bassa (una soglia alta sarebbe
        irraggiungibile). Il floor evita gli spawn gi\`a dentro il raggio di successo, tranne
        quando l'intero plume \`e pi\`u piccolo del floor stesso, nel qual caso il disco si
        restringe alla sua estensione;
  \item se $d_{\max}-d_{\min} < 500\,\text{m}$ si abbassa $d_{\min}$ (fino al floor) per
        garantire un intervallo abbastanza ampio da non far ripartire tutti gli agenti dalla
        stessa distanza.
\end{itemize}
Si estrae quindi una distanza obiettivo $d^\star \sim \mathcal{U}(d_{\min},\,d_{\max})$ e si
sceglie a caso una cella valida entro $\pm 50\,\text{m}$ da $d^\star$ (con piccolo jitter),
ripetendo finch\'e i vincoli di dominio, concentrazione e distanza dalla costa sono
soddisfatti. Campionando in modo uniforme $d^\star$, gli spawn si \textbf{spalmano su tutto}
$[d_{\min},\,d_{\max}]$ invece di concentrarsi sull'anello esterno.

\begin{lstlisting}
d = dist_from_source(plume_cells)          # celle con c>=0.5 e land>=50 m
d_disk = d[d <= 3000]                       # tetto d_cap = 3 km
d_max  = min(3000, d_disk.max())
d_min  = max(300, percentile(d_disk, 50))   # mediana, floor 300 m
if d_max - d_min < 500:                      # ampiezza minima dell'intervallo
    d_min = max(300, d_max - 500)
target = uniform(d_min, d_max)               # UNIFORME in distanza
band   = plume_cells[abs(d_disk - target) <= 50]   # cella entro +-50 m
return random.choice(band)                   # (con jitter +-5 m)
\end{lstlisting}

\begin{figure}[H]
\centering
\begin{tikzpicture}[font=\footnotesize]
  % 1 unita' = 500 m; disco esterno d_max, foro interno d_min
  \fill[blue!10] (0,0) circle (4.2);
  \fill[white]   (0,0) circle (1.5);
  \draw[dashed, thick, blue!55!black] (0,0) circle (1.5);
  \draw[dashed, thick, blue!55!black] (0,0) circle (4.2);
  % spawn estratti uniformemente in distanza nell'anello [d_min, d_max]
  \foreach \a/\r in {18/1.8, 52/3.5, 88/2.4, 124/3.9, 158/2.0, 196/3.2, 228/2.7,
                     262/4.0, 298/1.7, 332/3.0, 8/3.6, 72/2.9, 140/1.6, 210/2.3,
                     280/3.4, 350/2.6}{
    \fill[orange!85!black] (\a:\r) circle (1.6pt);
  }
  \node[star, star points=5, fill=yellow, draw, minimum size=5mm, inner sep=0pt] at (0,0) {};
  \draw[-{Latex[length=2mm]}, gray!65] (0,0) -- (68:1.5)
       node[pos=0.6, sloped, above, font=\scriptsize] {$d_{\min}$};
  \draw[-{Latex[length=2mm]}, gray!65] (0,0) -- (255:4.2)
       node[pos=0.72, sloped, below, font=\scriptsize] {$d_{\max}$};
\end{tikzpicture}
\caption{Disco di spawn adattivo $[d_{\min},d_{\max}]$ (versione adottata). $d_{\max}$ \`e il
         tetto di $3\,\text{km}$ ridotto all'estensione del plume; $d_{\min}$ \`e la mediana
         delle distanze delle celle di plume (\emph{floor} $300\,\text{m}$) e scala con la
         diffusione. I punti di partenza (arancioni) sono estratti in modo \emph{uniforme
         rispetto alla distanza} nell'anello, coprendo l'intero intervallo anzich\'e
         addensarsi sul bordo esterno.}
\label{fig:spawn_disk}
\end{figure}

\textbf{Pro.} Il disco \`e \emph{adattivo} (si ridimensiona con l'estensione del plume) e
distribuisce gli spawn in modo uniforme su tutto $[d_{\min},d_{\max}]$: la distanza iniziale
copre l'intero intervallo, rendendo il confronto fra modelli pi\`u equo e statisticamente
pi\`u solido. Il floor a $300\,\text{m}$ e l'ampiezza minima mantengono comunque il compito
non banale, senza inizializzazioni gi\`a prossime al bersaglio.
\textbf{Contro.} La logica \`e pi\`u complessa (percentile, banda di tolleranza, fallback),
e la distanza tipica di partenza non \`e un valore fisso ma dipende dalla mediana delle celle
di plume del singolo scenario, quindi varia da caso a caso. Soprattutto, campionando in modo
uniforme su $[d_{\min},d_{\max}]$, l'agente pu\`o partire anche da \textbf{distanze brevi}
(prossime a $d_{\min}$), mentre la cascata sceglieva sempre il punto valido pi\`u lontano:
alcuni episodi risultano quindi mediamente pi\`u facili rispetto al primo algoritmo.

\paragraph{Perch\'e il passaggio, e quale algoritmo usa cosa.}
Con la velocit\`a fissa la cascata era adeguata: il compito era gi\`a difficile e la distanza
di partenza, pur concentrata, risultava paragonabile fra episodi. Introducendo la
\textbf{velocit\`a variabile} (Sezione~\ref{sec:rl_velocity}) il confronto avviene per\`o
\emph{fra} modelli e \emph{fra} velocit\`a massime diverse: una distribuzione di spawn
degenere avrebbe reso i confronti sensibili al particolare anello estratto pi\`u che alla
reale capacit\`a dell'agente. Il disco adattivo, uniforme in distanza e indipendente dalla
forma istantanea del plume, offre una base di valutazione controllata e riproducibile. In
pratica: i modelli a \textbf{velocit\`a fissa} dei primi esperimenti usano la
\textbf{cascata}; tutti i modelli e le analisi a \textbf{velocit\`a variabile} --- l'intero
range di velocit\`a (Sezione~\ref{sec:rl_velocity}), i $v_{\max}$ intermedi, il confronto fra
corona singola e doppia e il test a singolo agente senza formazione --- usano il
\textbf{disco adattivo}.

L'analisi delle \emph{spawn map} (Sezione~\ref{sec:spawn_map}) verifica infine, per il disco,
che l'esito di un episodio non dipenda dalla particolare posizione di partenza.

\section{Il segnale di reward}
\label{sec:reward}

Oltre a osservazione e dinamica, l'ambiente calcola ad ogni passo un \textbf{segnale di
reward} scalare $r_t$. Esso \`e usato \emph{unicamente} dall'agente RL, che ne massimizza
il ritorno atteso in fase di addestramento (Capitolo~\ref{cap:rl}); il metodo del gradiente
lo ignora, poich\'e decide l'azione analiticamente dalle sole letture correnti. Nella sua
forma finale, adottata come riferimento, il reward combina un piccolo insieme di segnali
essenziali:
\begin{equation}
    r_t = \underbrace{r^{\text{reach}}_t}_{+100\ \text{a successo}}
        \;+\; \underbrace{r^{\text{dist}}_t}_{\text{avvicinamento}}
        \;+\; \underbrace{r^{\text{grad}}_t}_{\pm 0{,}05}
        \;+\; \underbrace{r^{\text{time}}_t}_{-0{,}1/\text{step}}
        \;+\; \underbrace{r^{\text{land}}_t}_{\text{costa}}
        \;+\; \underbrace{r^{\text{bound}}_t}_{-10\ \text{fuori dominio}}
    \label{eq:reward}
\end{equation}
dove: $r^{\text{reach}}_t = +100$ (pi\`u un piccolo bonus per l'arrivo anticipato) premia il
raggiungimento della sorgente; $r^{\text{dist}}_t$
\`e proporzionale alla \emph{riduzione} di distanza dalla sorgente nel passo (reward di
distanza lineare); $r^{\text{grad}}_t = \pm 0{,}05$ premia il muoversi lungo il gradiente
crescente di concentrazione; $r^{\text{time}}_t = -0{,}1$ per step incentiva a raggiungere
la sorgente rapidamente; $r^{\text{land}}_t$ penalizza la prossimit\`a eccessiva alla costa
(via la mappa EDT); $r^{\text{bound}}_t = -10$ penalizza l'uscita dal dominio. In codice, il
reward si assembla ad ogni passo come somma di questi contributi:

\begin{lstlisting}
def _compute_reward(self):
    r = 0.0
    if self._source_reached():                    # 1) raggiungimento (+ bonus di tempo)
        r += 100.0 + max(0, (N_max - k) / N_max * 50)
    if self._out_of_bounds():                     # 2) uscita dal dominio
        r += -10.0
    r += (prev_dist - curr_dist) * K_dist          # 3) distanza (lineare)
    r += 0.05 if conc > prev_conc else -0.05       # 4) gradiente di concentrazione
    r += -0.1                                       # 5) penalita' di tempo (per step)
    r += land_proximity_penalty(land_dist)         # 6) prossimita' alla costa (EDT)
    return r
\end{lstlisting}

La scelta di questa forma \emph{minimale} non \`e a priori: \`e il risultato di uno studio
di ablation (Sezione~\ref{sec:rl_reward}) che ha confrontato varianti pi\`u ricche ---
con reward di allineamento a vento e corrente, distanza quadratica, \emph{zone multiplier},
\emph{retreat asymmetry} e penalit\`a di stagnazione --- mostrando che i segnali aggiuntivi
non migliorano, e in alcuni casi peggiorano, le prestazioni. Il gradiente di concentrazione,
calcolato direttamente sul campo simulato, incorpora gi\`a implicitamente l'effetto di vento
e corrente.

